\documentclass[11pt, a4paper]{article}

% ─── Packages ─────────────────────────────────────────────────────────
\usepackage[utf8]{inputenc}
\usepackage[T1]{fontenc}
\usepackage{lmodern}
\usepackage[margin=1in]{geometry}
\usepackage{hyperref}
\usepackage{xcolor}
\usepackage{listings}
\usepackage{booktabs}
\usepackage{longtable}
\usepackage{array}
\usepackage{tabularx}
\usepackage{enumitem}
\usepackage{titlesec}
\usepackage{fancyhdr}
\usepackage{microtype}
\usepackage{graphicx}
\usepackage{textcomp}

% ─── Colours ──────────────────────────────────────────────────────────
\definecolor{rsnavy}{HTML}{0B1424}
\definecolor{rsblue}{HTML}{1D4ED8}
\definecolor{rsred}{HTML}{DC2626}
\definecolor{rsteal}{HTML}{0F766E}
\definecolor{rsgrey}{HTML}{475569}
\definecolor{rscode}{HTML}{F1F5F9}

% ─── Hyperlinks ───────────────────────────────────────────────────────
\hypersetup{
  colorlinks=true,
  linkcolor=rsblue,
  urlcolor=rsblue,
  citecolor=rsblue,
  pdftitle={RoadSOS Technical Documentation},
  pdfauthor={RoadSOS Engineering},
  pdfsubject={Architecture and Implementation Reference}
}

% ─── Listings ─────────────────────────────────────────────────────────
\lstdefinestyle{rsstyle}{
  backgroundcolor=\color{rscode},
  basicstyle=\ttfamily\footnotesize,
  breakatwhitespace=false,
  breaklines=true,
  captionpos=b,
  keepspaces=true,
  numbers=left,
  numbersep=6pt,
  numberstyle=\tiny\color{rsgrey},
  showspaces=false,
  showstringspaces=false,
  showtabs=false,
  tabsize=2,
  frame=single,
  framesep=4pt,
  rulecolor=\color{rsgrey!30},
  xleftmargin=14pt
}
\lstset{style=rsstyle}

% ─── Headings ─────────────────────────────────────────────────────────
\titleformat{\section}{\Large\bfseries\color{rsnavy}}{\thesection}{0.6em}{}
\titleformat{\subsection}{\large\bfseries\color{rsblue}}{\thesubsection}{0.6em}{}
\titleformat{\subsubsection}{\normalsize\bfseries\color{rsnavy}}{\thesubsubsection}{0.6em}{}

% ─── Header / Footer ──────────────────────────────────────────────────
\pagestyle{fancy}
\fancyhf{}
\fancyhead[L]{\small\color{rsgrey} RoadSOS Technical Documentation}
\fancyhead[R]{\small\color{rsgrey} \nouppercase{\leftmark}}
\fancyfoot[C]{\small\thepage}
\renewcommand{\headrulewidth}{0.4pt}

% ─── Custom command for inline file references ────────────────────────
\newcommand{\file}[1]{\texttt{\small\color{rsblue}#1}}
\newcommand{\code}[1]{\texttt{\small#1}}

% ─── Title ────────────────────────────────────────────────────────────
\title{
  \vspace{-1.5cm}
  \Huge\bfseries\color{rsnavy} RoadSOS \\
  \large\normalfont\color{rsgrey} Technical Documentation \\[0.4em]
  \normalsize\color{rsgrey} Architecture, Algorithms, and Implementation Reference
}
\author{
  \small\color{rsgrey} \texttt{github.com/Arthrevs/Roadproj} \\
  \small\color{rsgrey} Branch: \texttt{main} \quad Commit: \texttt{f94dc2a}
}
\date{\small\color{rsgrey} May 2026}

% ──────────────────────────────────────────────────────────────────────
% DOCUMENT
% ──────────────────────────────────────────────────────────────────────
\begin{document}

\maketitle
\thispagestyle{empty}
\vspace{0.5cm}

\begin{center}
\colorbox{rscode}{\parbox{0.85\textwidth}{\centering\small
\textbf{Scope.} This document is the engineering reference for RoadSOS. It covers the request lifecycle, the OpenStreetMap Overpass query layer, the Google Places fallback strategy, the four-tier offline cache hierarchy, the AI triage model, internationalization across 43 languages, the real-time map renderer, crash detection signals, and the deployment topology. Every number cited is grounded in source. Every external dependency is named. No marketing.
}}
\end{center}

\vspace{0.4cm}
{\small\tableofcontents}
\newpage

% ──────────────────────────────────────────────────────────────────────
\section{System Overview}
% ──────────────────────────────────────────────────────────────────────

RoadSOS is a Progressive Web Application that locates emergency-services contacts (hospitals, police, ambulances, fire stations, towing, repair, tyre/puncture, showrooms) within driving distance of a user's GPS position and dispatches their coordinates to pre-configured emergency contacts via WhatsApp or SMS.

\subsection{Top-Level Topology}

\begin{lstlisting}[language={},numbers=none,frame=none,backgroundcolor=\color{white}]
+---------------------------------------------------------------+
|                   BROWSER  (PWA, installable)                 |
|                                                               |
|   React 18  +  Vite 8  +  Workbox 7 (Service Worker)          |
|   |- Leaflet 1.9.4    (real OSM map, dark tiles)              |
|   |- i18next 26       (43 languages, RTL support)             |
|   |- localStorage     (24h search-result cache)               |
|   `- bundled JSON     (249 facilities x 196 countries)        |
+----------------------------+----------------------------------+
                             |  HTTPS  (CORS-allowlisted)
                             v
+---------------------------------------------------------------+
|              FASTAPI BACKEND  (Render, Python 3.11.9)         |
|                                                               |
|   uvicorn + asyncio + httpx                                   |
|   |- /search    (4-phase orchestrator)                        |
|   |- /triage    (Gemini 2.0 Flash + rule-based fallback)      |
|   |- /dispatch  (telemetry sink for SOS events)               |
|   |- /health    (liveness + version)                          |
|   `- /offline-pack  (static facility bundle for warm cache)   |
+----+--------------------+--------------------+----------------+
     |                    |                    |
     v                    v                    v
+----------+        +------------+      +--------------+
|   OSM    |        |   Google   |      |  Nominatim   |
| Overpass |        |   Places   |      | (geocoder)   |
|  (free)  |        |   (paid)   |      |   (free)     |
+----------+        +------------+      +--------------+
\end{lstlisting}

\subsection{Tech Stack at a Glance}

\begin{tabularx}{\textwidth}{l X}
\toprule
\textbf{Layer} & \textbf{Components} \\
\midrule
Frontend runtime & React 18.3.1, React DOM 18.3.1, React Router 7.15.0 \\
Build tool & Vite 8.0.12, vite-plugin-pwa 1.3.0 (\code{injectManifest}) \\
Service worker & Workbox 7.4.1 (precaching, expiration, routing, strategies) \\
Map & Leaflet 1.9.4 + react-leaflet 4.2.1; CartoDB Dark Matter tiles \\
i18n & i18next 26.2.0 + react-i18next 17.0.8; 43 locales \\
Icons & lucide-react 1.16.0 \\
Backend runtime & FastAPI 0.115.0 + uvicorn 0.30.6, Python 3.11.9 \\
Async HTTP & httpx 0.27.2 \\
AI & Google Gemini 2.0 Flash via REST (no SDK; uses httpx) \\
Phone normalization & phonenumbers 8.13.46 \\
Storage & SQLite via aiosqlite 0.20.0 (offline-pack only) \\
Validation & pydantic 2.9.2 \\
Lint / format / test & ruff 0.11.13, pytest 8.3.3, pytest-asyncio 0.24.0, vitest 2.1.4 \\
Deployment & Render (backend), Vercel (frontend) \\
CI/CD & GitHub Actions: \code{frontend-ci.yml}, \code{backend-tests.yml}, \code{pr-guard.yml} \\
\bottomrule
\end{tabularx}

% ──────────────────────────────────────────────────────────────────────
\section{Backend: The \texttt{/search} Orchestrator}
% ──────────────────────────────────────────────────────────────────────

The search endpoint resolves the question \emph{``what emergency services are near these coordinates, and what are their phone numbers?''} It is implemented in \file{backend/services/search\_service.py} as a four-phase asynchronous pipeline.

\subsection{Phase 1 --- Parallel Geocode and Overpass}

Two operations begin simultaneously via \code{asyncio.gather}: a reverse-geocode against Nominatim and an Overpass query against OpenStreetMap.

\begin{lstlisting}[language=Python,caption={search\_service.py --- Phase 1 (lines 127--129)}]
geo_task = asyncio.create_task(_safe_geocode(lat, lon))
osm_task = asyncio.create_task(_safe_overpass(lat, lon))
geo, osm_contacts = await asyncio.gather(geo_task, osm_task)
\end{lstlisting}

Wall-clock latency is therefore $\max(t_{\text{geocode}}, t_{\text{overpass}})$ rather than the sum. Both are wrapped in \code{\_safe\_*} helpers that catch every exception and return a valid fallback shape; failures never propagate.

\subsubsection{The Overpass Query Layer}

\file{backend/services/overpass\_service.py} builds an Overpass QL query that selects nine categories within a configurable radius (default 5{,}000 m). The query is structured as a union over tag filters:

\begin{lstlisting}[language={},caption={Overpass QL --- categories targeted (overpass\_service.py:85--119)}]
[out:json][timeout:25];
(
  nwr["amenity"="hospital"](around:5000, LAT, LON);
  nwr["amenity"="clinic"] ["emergency"="yes"](around:5000, LAT, LON);
  nwr["amenity"="police"](around:5000, LAT, LON);
  nwr["amenity"="fire_station"](around:5000, LAT, LON);
  nwr["amenity"="car_repair"](around:5000, LAT, LON);
  nwr["shop"="car_repair"](around:5000, LAT, LON);
  nwr["shop"="tyres"](around:5000, LAT, LON);
  nwr["shop"="car"](around:5000, LAT, LON);
  nwr["emergency"="ambulance_station"](around:5000, LAT, LON);
);
out center;
\end{lstlisting}

If the first response yields fewer than three results, the radius doubles to 10{,}000 m and the query repeats. This handles sparse regions without paying for the wider radius on dense urban queries.

\subsubsection{Multi-Mirror Retry Strategy}

A single Overpass endpoint frequently rate-limits or 502s. The fetcher rotates across three independent mirrors with exponential backoff:

\begin{tabularx}{\textwidth}{l X}
\toprule
\textbf{Attempt} & \textbf{Endpoint and Backoff} \\
\midrule
1 & \code{https://overpass-api.de/api/interpreter} (no delay) \\
2 & \code{https://overpass.kumi.systems/api/interpreter} (2 s) \\
3 & \code{https://overpass.openstreetmap.fr/api/interpreter} (4 s) \\
\bottomrule
\end{tabularx}

Each request has a 30-second timeout. If all three mirrors fail, the function returns \code{[]} rather than raising; downstream phases continue with whatever data exists.

\subsubsection{Element Parsing and Distance}

Each Overpass element is normalised by \code{parse\_element} (overpass\_service.py:122--170):

\begin{enumerate}[leftmargin=*, itemsep=2pt, topsep=2pt]
  \item Extract \code{lat}, \code{lon} (from the element or its \code{center} field for ways).
  \item Compute haversine distance from the user, expressed in kilometres rounded to two decimals.
  \item Map the OSM tag to one of nine internal categories: \texttt{hospital}, \texttt{police}, \texttt{ambulance}, \texttt{fire}, \texttt{towing}, \texttt{repair}, \texttt{tyre}, \texttt{showroom}.
  \item Drop results with no name \emph{and} no phone (reduces ``Unnamed Building'' noise).
  \item Parse the \code{opening\_hours} tag into a boolean \code{isOpen} via \file{services/hours\_parser.py} (supports \code{24/7}, day ranges, time ranges, and holiday tags).
\end{enumerate}

The haversine implementation is unrolled in pure Python (no GIS dependency):

\begin{lstlisting}[language=Python,caption={overpass\_service.py:66--78}]
def haversine(lat1, lon1, lat2, lon2):
    R = 6371.0  # Earth radius (km)
    dlat = math.radians(lat2 - lat1)
    dlon = math.radians(lon2 - lon1)
    a = (math.sin(dlat / 2) ** 2 +
         math.cos(math.radians(lat1)) *
         math.cos(math.radians(lat2)) *
         math.sin(dlon / 2) ** 2)
    return 2 * R * math.asin(math.sqrt(a))
\end{lstlisting}

\subsubsection{Proximity Deduplication}

OSM frequently lists the same facility as both a node (the entrance) and a way (the building polygon). \code{\_dedupe\_smart} (overpass\_service.py:185--210) keeps results that are either more than 50 m apart \emph{or} have substantially different names; otherwise the earlier-seen entry wins.

\subsection{Phase 2 --- Conditional Google Places Fallback}

Phase 2 fires only if Phase 1 produced fewer than three contacts with a dialable phone number. This conserves quota: in dense OSM regions (most of Europe, India, and East Asia), Google is never called.

\begin{lstlisting}[language=Python,caption={search\_service.py --- Phase 2 trigger (lines 135--138)}]
phoned_osm = [c for c in osm_contacts if c.get("phone")]
google_contacts: list[dict] = []
if len(phoned_osm) < 3:
    google_contacts = await _safe_google(lat, lon, geo.get("country_code"))
\end{lstlisting}

\file{services/googleplaces\_service.py} issues four parallel Nearby Search calls (hospital, police, car\_repair, fire\_station) each capped at five results. Per-result Place Details lookups add the \code{formatted\_phone\_number} field. The reverse-geocoded country code from Phase 1 is passed as a region hint to improve relevance.

\subsubsection{Multi-Key Rotation}

The \code{Mapsplatformkey} environment variable is parsed as a comma-separated list of API keys. Round-robin rotation distributes load across multiple billing accounts:

\begin{lstlisting}[language=Python,caption={googleplaces\_service.py:30--36}]
def _load_keys() -> list[str]:
    multi = os.getenv("Mapsplatformkey", "")
    keys = [k.strip() for k in multi.split(",") if k.strip()]
    if not keys:
        single = os.getenv("Mapsplatformkey", "")
        keys = [single] if single else []
    return keys
\end{lstlisting}

If no key is configured, Google is silently disabled --- the search still returns whatever Overpass produced.

\subsection{Phase 3 --- Merge and Deduplicate}

The two contact lists are merged and deduplicated by phone digits first, then by lowercased name:

\begin{lstlisting}[language=Python,caption={search\_service.py:42--64 --- deduplication}]
def deduplicate(contacts: list[dict]) -> list[dict]:
    out, seen_names = [], set()
    for c in contacts:
        if not isinstance(c, dict): continue
        name_key = (c.get("name") or "").lower().strip()
        if not name_key or name_key in seen_names: continue
        if any(phones_match(c.get("phone"), e.get("phone")) for e in out):
            continue
        seen_names.add(name_key)
        out.append(c)
    return out
\end{lstlisting}

\code{phones\_match} (phone\_utils.py) compares the last ten digits after \code{phonenumbers}-normalisation. This handles cases where OSM and Google disagree on transliteration (``GS Custom'' vs.\ ``gs sustom'') but the phone number is canonical. The merged list is then sorted by distance.

\subsection{Phase 4 --- Phone Enrichment}

For the top six closest results without a phone number, the system performs a Find Place from Text lookup against Google Places, biased to a 500-metre radius. On miss, a fallback Nearby Search within 100 m provides a second chance. Each enrichment costs at most one Place Details billable call, capped at six per search to bound worst-case spend.

\subsection{Response Shape}

\begin{lstlisting}[language=json,caption={Sample /search response}]
{
  "contacts": [
    {
      "id": "node/8127364521",
      "name": "Apollo Hospitals, Bannerghatta",
      "category": "hospital",
      "phone": "080-26793000",
      "lat": 12.8932, "lon": 77.5972,
      "distance": 1.42,
      "source": "OpenStreetMap",
      "isOpen": true
    }
  ],
  "source": "OpenStreetMap + Google Places",
  "landmark": "Bannerghatta Road, BTM Layout, Bengaluru",
  "country_code": "IN",
  "count": 14
}
\end{lstlisting}

% ──────────────────────────────────────────────────────────────────────
\section{Backend: Caching}
% ──────────────────────────────────────────────────────────────────────

\file{backend/services/cache.py} implements an async-safe in-memory TTL cache with LRU eviction. Three module singletons are configured for the dominant lookups:

\begin{tabularx}{\textwidth}{l c c X}
\toprule
\textbf{Cache} & \textbf{TTL} & \textbf{Max entries} & \textbf{Key strategy} \\
\midrule
\code{overpass\_cache} & 1 hour & 200 & \code{round(lat,4)\_round(lon,4)\_r5000} \\
\code{google\_cache} & 1 hour & 200 & coordinate key + region code \\
\code{geocode\_cache} & 24 hours & 500 & coordinate key only \\
\bottomrule
\end{tabularx}

\subsection{Key Granularity}

Coordinates are rounded to four decimal places before keying. At the equator one degree of latitude is approximately 111 km, so $10^{-4}$ degrees $\approx$ 11 metres. Two users standing 10 m apart resolve to the same cache key and share results --- semantically correct for ``what is near me'' queries, and trivially exploitable for warm-cache benchmarks (a repeated query from the same demo location returns in tens of milliseconds).

\subsection{Eviction and Concurrency}

\begin{itemize}[leftmargin=*, itemsep=2pt, topsep=2pt]
  \item \textbf{LRU on capacity.} When \code{len(cache) >= max\_entries}, the entry with the oldest timestamp is removed.
  \item \textbf{Opportunistic expiry.} Each \code{set} call scans for and removes expired entries before inserting.
  \item \textbf{Async lock.} Each cache instance holds an \code{asyncio.Lock}; concurrent \code{get}/\code{set} operations are serialised per cache.
\end{itemize}

\subsection{Negative-Result Policy}

Overpass results with zero phone numbers are not cached. This is deliberate: a phoneless cache hit would short-circuit Phase 2 (Google enrichment), permanently denying the user phone-equipped contacts they could otherwise get. The cost is recomputation for sparse regions, accepted as a quality-over-latency trade.

% ──────────────────────────────────────────────────────────────────────
\section{Backend: Rate Limiting}
% ──────────────────────────────────────────────────────────────────────

\file{services/rate\_limiter.py} implements a per-IP token bucket. Each client receives a bucket of capacity 10 that refills at the configured rate; requests consume one token. When the bucket is empty, the request receives \code{429 Too Many Requests} with a \code{Retry-After} header in seconds.

\begin{tabularx}{\textwidth}{l c c X}
\toprule
\textbf{Endpoint} & \textbf{Rate} & \textbf{Burst} & \textbf{Notes} \\
\midrule
\code{GET /search} & 30 req/min & 10 & Per-IP \\
\code{POST /triage} & 20 req/min & 10 & Per-IP, lower because AI calls are costlier \\
\bottomrule
\end{tabularx}

\subsection{Render Reverse-Proxy Awareness}

Render terminates TLS at its edge and forwards requests with an \code{X-Forwarded-For} header. The default ASGI \code{request.client.host} reports the Render edge IP, which would collapse all users into a single bucket. \code{get\_client\_ip} (rate\_limiter.py:70) prefers the first entry of \code{X-Forwarded-For} when present:

\begin{lstlisting}[language=Python,numbers=none,frame=none]
xff = request.headers.get("x-forwarded-for", "")
return xff.split(",")[0].strip() if xff else request.client.host
\end{lstlisting}

Buckets are stored in a plain Python dictionary protected by \code{asyncio.Lock}. They reset whenever the Render dyno cold-starts (free tier idle timeout: 15 minutes).

% ──────────────────────────────────────────────────────────────────────
\section{Backend: AI Triage}
% ──────────────────────────────────────────────────────────────────────

\file{backend/services/ai\_triage.py} reorders a contact list to match the user's situation. The model is \texttt{gemini-2.0-flash}, called via direct REST against \url{https://generativelanguage.googleapis.com/v1beta/models/} (no SDK dependency --- uses the same \code{httpx} client as the rest of the backend). The free tier allows 60 requests per minute and 1500 requests per day with no billing requirement. Two boolean situational flags are accepted: \code{injured} and \code{blocking\_traffic}. JSON mode is opt-in via \code{generationConfig.responseMimeType = "application/json"}.

\subsection{Prompt Contract}

The system prompt mandates a strict JSON response with exactly two top-level keys:

\begin{lstlisting}[language={},caption={Required model output schema}]
{
  "contacts": [<same items as input, reordered>],
  "reason": "<one short sentence explaining the top pick>"
}
\end{lstlisting}

Validation checks (ai\_triage.py:87--99) that:

\begin{enumerate}[leftmargin=*, itemsep=2pt, topsep=2pt]
  \item the response parses as JSON,
  \item \code{contacts} is a list of exactly the same length as the input,
  \item every contact carries the same \code{id} as an input entry,
  \item \code{reason} is a non-empty string.
\end{enumerate}

\subsection{Rule-Based Fallback}

If the model call fails (no API key, network error, malformed response, validation failure, count mismatch), a deterministic priority table takes over:

\begin{tabularx}{\textwidth}{l X}
\toprule
\textbf{Situation} & \textbf{Priority order (descending)} \\
\midrule
Injured \& blocking & ambulance, hospital, police, towing, repair \\
Injured only & ambulance, hospital, police, repair, towing \\
Blocking only & police, towing, ambulance, hospital, repair \\
Neither & repair, tyre, police, towing, hospital, ambulance \\
\bottomrule
\end{tabularx}

The fallback is the same logic the model is prompted to follow, so the user-visible behaviour is consistent whether or not the model responded.

\subsection{Independence from Search}

Triage is a separate endpoint (\code{POST /triage}) called only after the user answers the situational questions. The map and contact list render with no AI dependency; if Gemini is entirely unreachable, the list is still useful.

% ──────────────────────────────────────────────────────────────────────
\section{Backend: Reliability Guarantees}
% ──────────────────────────────────────────────────────────────────────

Every external call is wrapped in a \code{\_safe\_*} helper that catches all exceptions and returns a valid fallback shape:

\begin{tabularx}{\textwidth}{l X}
\toprule
\textbf{Wrapper} & \textbf{Fallback return} \\
\midrule
\code{\_safe\_overpass} & \code{[]} \\
\code{\_safe\_google} & \code{[]} \\
\code{\_safe\_geocode} & \code{\{"landmark": "lat°, lon°", "country\_code": null\}} \\
\code{deduplicate} & Skips non-dict items rather than raising \\
\code{triage\_service.prioritise} & Rule-based ordering with outer try/except as last resort \\
\bottomrule
\end{tabularx}

\subsection{Catastrophic Failure Path}

If every upstream is unavailable simultaneously:

\begin{lstlisting}[language=json,numbers=none,caption={Response when all upstreams fail}]
{
  "contacts": [],
  "source": "none (upstreams unavailable)",
  "landmark": "13.0827, 80.2707",
  "country_code": null,
  "count": 0
}
\end{lstlisting}

The HTTP status is still \texttt{200 OK}. The frontend interprets an empty result as a trigger for the offline fallback chain (Section~\ref{sec:offline}).

\subsection{Middleware Safety Net}

\file{backend/middleware.py} adds three layers around every request:

\begin{enumerate}[leftmargin=*, itemsep=2pt, topsep=2pt]
  \item \textbf{RequestIDMiddleware.} Stamps every request with a UUID-v4 (or echoes the client's \code{X-Request-ID}) and exposes it on the response.
  \item \textbf{RequestLogMiddleware.} Emits one log line per request: \code{method path status duration\_ms request\_id[:8]}.
  \item \textbf{ErrorHandlingMiddleware.} Outermost. Catches anything an inner handler raised, logs the full traceback at \texttt{ERROR} with the request ID, and returns \code{503 Service Unavailable} with the request ID in the JSON body for support reporting.
\end{enumerate}

% ──────────────────────────────────────────────────────────────────────
\section{Frontend: Offline Architecture}
\label{sec:offline}
% ──────────────────────────────────────────────────────────────────────

The frontend implements a four-tier fallback chain. Each tier is queried only if the preceding tier has no answer. The order is fixed in \file{frontend/src/App.jsx} (lines 197--255).

\subsection{Tier 1 --- Live Backend Search}

\code{searchNearby(lat, lon)} in \file{src/utils/overpass.js} posts to \code{/search?lat=\&lon=} on the configured backend. Timeout: 8 seconds. On success the response is also written to Tier 2 for future offline use.

\subsection{Tier 2 --- localStorage Cache}

\file{src/utils/offlineDB.js} stores responses in \code{window.localStorage} keyed by a coarsened coordinate grid:

\begin{lstlisting}[language=JavaScript,caption={offlineDB.js: 24-hour client-side cache}]
const TTL_MS = 24 * 60 * 60 * 1000; // 24 hours

function locationKey(lat, lon) {
  // Round to 2 decimal places (~1.1 km grid)
  return `roadsos_cache_${lat.toFixed(2)}_${lon.toFixed(2)}`;
}
\end{lstlisting}

The grid is 1.1 km square at the equator. A user who searched yesterday from any point within the same square gets a cache hit today. Entries older than the TTL are removed on read.

\subsection{Tier 3 --- Bundled Facilities}

When network requests fail and no cache entry exists, \file{src/utils/bundledFacilities.js} computes the nearest entries from a JSON file shipped in the build bundle: \file{src/data/bundled\_facilities.json} (249 records, 196 countries).

\begin{tabularx}{\textwidth}{l c X}
\toprule
\textbf{Region} & \textbf{Records} & \textbf{Notes} \\
\midrule
India (deepest) & $\sim$40 & Tamil Nadu + Delhi/Mumbai/Bengaluru/Hyderabad metros \\
North America & 22 & US trauma centers (Level I + II), Canadian provinces \\
Europe & 38 & EU member trauma units + UK NHS major centres \\
East and Southeast Asia & 26 & Japan, China, Korea, Thailand, Vietnam, Indonesia, Philippines, Singapore, Malaysia \\
Middle East and North Africa & 18 & UAE, Saudi Arabia, Israel, Turkey, Iran, Egypt \\
Sub-Saharan Africa & 32 & At least one per UN-recognised state \\
Latin America \& Caribbean & 24 & Brazil/Argentina/Mexico majors plus smaller-nation hospitals \\
Oceania & 8 & Australia, New Zealand, Pacific island states \\
\bottomrule
\end{tabularx}

Search radius starts at 80 km and expands to 600 km if zero results match. Distance is computed by the same haversine formula used server-side, ensuring frontend and backend agree on ordering.

\subsection{Tier 4 --- Mock Data}

If all three tiers above return empty, a hardcoded array of seven Bengaluru-area contacts (App.jsx:41--112) is used as a last visual placeholder. This tier exists so the UI is never blank; in production deployments it would be removed under a build flag.

\subsection{Service Worker}

\file{frontend/public/sw.js} is registered through \code{vite-plugin-pwa} with the \code{injectManifest} strategy. The Workbox configuration includes:

\begin{itemize}[leftmargin=*, itemsep=2pt, topsep=2pt]
  \item \textbf{Precache.} The full app shell (HTML, JS, CSS, fonts) is injected via \code{self.\_\_WB\_MANIFEST} at build time. Total: roughly 670 KiB across six entries.
  \item \textbf{Runtime cache for /search.} \code{NetworkFirst} with an 8-second timeout, 50-entry LRU, 24-hour max age. After the timeout, the SW serves the most recent cached response for that URL.
  \item \textbf{Runtime cache for /offline-pack.} \code{CacheFirst} with a single entry and seven-day max age. The backend's offline pack is downloaded once and reused.
  \item \textbf{skipWaiting() on install.} Critical for live deployments. Without it, a new build would only take effect after the user closed every open tab.
\end{itemize}

\subsection{Country Emergency Numbers}

Separate from facility lookup, \file{src/utils/emergencyNumbers.js} is a compile-time mapping from ISO-3166 alpha-2 country code to a list of national emergency dial codes. Examples: \code{IN} $\rightarrow$ \code{[108, 100, 101, 112]}, \code{US} $\rightarrow$ \code{[911]}, \code{GB} $\rightarrow$ \code{[999, 112]}, \code{DE} $\rightarrow$ \code{[112, 110]}. The banner at the top of the app renders these in the user's language and is fully offline-capable.

% ──────────────────────────────────────────────────────────────────────
\section{Frontend: The Real-Time Map}
% ──────────────────────────────────────────────────────────────────────

\file{frontend/src/components/RealMap.jsx} renders an OpenStreetMap-backed Leaflet map centred on the user's GPS coordinates.

\subsection{Tile Provider}

Tiles are fetched from CartoDB's Dark Matter raster style:

\begin{lstlisting}[language={},numbers=none,frame=none]
https://{s}.basemaps.cartocdn.com/dark_all/{z}/{x}/{y}{r}.png
\end{lstlisting}

CartoDB serves these tiles free of charge for non-commercial use with attribution. No API key is required. The \code{\{s\}} placeholder distributes across the subdomains \code{a}--\code{d}. The \code{\{r\}} suffix delivers Retina (\code{@2x}) tiles on high-DPI displays.

\subsection{Attribution}

OSM's Open Database License (ODbL) and CartoDB's terms both require visible attribution on derived works. The Leaflet attribution control renders bottom-right:

\begin{lstlisting}[language=HTML,numbers=none,frame=none]
&copy; <a href="https://www.openstreetmap.org/copyright">OSM</a>
&copy; <a href="https://carto.com/attributions">CARTO</a>
\end{lstlisting}

\subsection{User Marker}

A custom \code{L.divIcon} draws three concentric layers:

\begin{enumerate}[leftmargin=*, itemsep=2pt, topsep=2pt]
  \item \textbf{Halo.} 110 px radial gradient at 33\% opacity, providing visual mass.
  \item \textbf{Pulse ring.} 28 px circle animated via \code{@keyframes rs-blip} (2.5 s ease-out infinite).
  \item \textbf{Dot.} 18 px solid circle with a 3 px white border and a coloured drop-shadow.
\end{enumerate}

When the GPS signal is older than the freshness threshold, the colour shifts from green (\code{\#22C55E}) to grey (\code{\#94A3B8}) and the pulse animation continues at reduced opacity.

\subsection{Service Markers}

Up to six contacts are rendered as map pins using their actual lat/lon coordinates from the search response. Each pin combines:

\begin{itemize}[leftmargin=*, itemsep=2pt, topsep=2pt]
  \item A coloured circle (red for medical: hospital, ambulance, fire; blue for police; teal for mechanical: towing, repair, tyre, showroom)
  \item A category emoji
  \item A floating chip above with the truncated facility name and distance in kilometres
\end{itemize}

The map is configured non-interactively in the hero section: \code{dragging}, \code{scrollWheelZoom}, \code{doubleClickZoom}, \code{touchZoom}, \code{boxZoom}, and \code{keyboard} are all disabled. This preserves the SOS button as the primary interaction. A separate full-screen mode (not yet shipped) would enable interaction.

\subsection{Re-Centering}

A child component, \code{MapRecenter}, hooks into \code{useMap()} and calls \code{map.setView()} with a 0.6-second animation whenever the location prop changes. This makes the demo location picker (Bengaluru $\rightarrow$ London $\rightarrow$ Tokyo) glide smoothly across the globe.

% ──────────────────────────────────────────────────────────────────────
\section{Frontend: Internationalization}
% ──────────────────────────────────────────────────────────────────────

\file{frontend/src/i18n/} contains 43 translation bundles registered through \code{i18next}.

\subsection{Language Coverage}

\subsubsection{Indian Languages (all 22 Schedule-VIII languages plus English)}

Assamese, Bengali, Bodo, Dogri, English, Gujarati, Hindi, Kannada, Kashmiri, Konkani, Maithili, Malayalam, Manipuri, Marathi, Nepali, Odia, Punjabi, Sanskrit, Santali, Sindhi, Tamil, Telugu, Urdu.

\subsubsection{Global Languages (most-spoken per region)}

Spanish, French, Portuguese, German, Italian, Dutch, Polish, Russian, Ukrainian, Romanian, Greek, Turkish, Arabic, Persian, Hebrew, Swahili, Amharic, Indonesian, Malay, Thai, Vietnamese, Chinese, Japanese, Korean, Filipino.

\subsection{Right-to-Left Support}

Six languages set the document direction to RTL on selection (Arabic, Hebrew, Persian, Urdu, Kashmiri, Sindhi). The handler in \file{src/i18n/index.js}:

\begin{lstlisting}[language=JavaScript,caption={i18n/index.js --- applyDir}]
function applyDir(code) {
  const loc = getLocale(code);
  document.documentElement.lang = loc.bcp47 || code;
  document.documentElement.dir = loc.dir || 'ltr';
}
\end{lstlisting}

\subsection{Selection Logic}

On first launch, the user sees a full-screen modal listing all 43 languages, grouped into ``India'' (22) and ``World'' (21) sections. There is no automatic pre-selection based on GPS country code; the user picks explicitly. The choice is persisted to \code{localStorage['roadsos:lang']} and never re-asked.

If the user dismisses the picker without choosing, the language defaults to:

\begin{enumerate}[leftmargin=*, itemsep=2pt, topsep=2pt]
  \item The previous session's choice from \code{localStorage}, if any.
  \item \code{navigator.language} (first two characters), if a translation bundle exists for it.
  \item \code{en}.
\end{enumerate}

\subsection{Translation Coverage}

Each bundle contains approximately 40 translation keys spanning SOS button states, emergency-category labels, location-status messages, action-button labels, crash-alert prompts, and Medical ID form fields. The complete key list is enumerated in \file{src/i18n/en.json} as the canonical reference.

% ──────────────────────────────────────────────────────────────────────
\section{Frontend: Crash Detection}
% ──────────────────────────────────────────────────────────────────────

\file{frontend/src/hooks/useLocation.js} fuses two sensor signals to identify a likely crash without requiring dedicated hardware.

\subsection{Signal 1 --- GPS Velocity Collapse}

The hook maintains a rolling buffer of GPS samples (each containing speed in m/s and a timestamp). A velocity collapse is registered when a sample at $\ge$ 25 km/h is followed within 2 seconds by a sample at $\le$ 5 km/h.

\subsection{Signal 2 --- Accelerometer Spike}

A \code{devicemotion} listener tracks \code{accelerationIncludingGravity}. A spike is registered when the magnitude exceeds 3.5 G (34.3 m/s$^2$). On iOS 13 and later, motion permission must be granted by an explicit user gesture; the first tap on the SOS area triggers \code{requestMotionPermission()}.

\subsection{Confirmation Logic}

A crash is reported only if both signals occur within 4 seconds of each other. The two-signal requirement substantially reduces false positives from phone drops (large G spike, no velocity change) and potholes (small velocity change, no sustained G spike).

\subsection{Cooldown}

After firing, the detector enters a 12-second cooldown during which neither signal can re-trigger the alert. This prevents an aftershock or a follow-up impact from doubling notifications.

% ──────────────────────────────────────────────────────────────────────
\section{Frontend: SOS Dispatch}
% ──────────────────────────────────────────────────────────────────────

\file{frontend/src/utils/sosDispatch.js} translates a button tap into a native messaging app deep link, choosing the dominant medium per country.

\subsection{WhatsApp-Dominant Countries}

For countries where WhatsApp is the de-facto messaging app, the handler issues \code{window.open('https://wa.me/<phone>?text=<encoded message>')}. The deep link opens WhatsApp (or its web client) pre-filled with:

\begin{itemize}[leftmargin=*, itemsep=2pt, topsep=2pt]
  \item The phrase ``Emergency. I need help.''
  \item GPS coordinates as decimal degrees.
  \item A Google Maps link with the same coordinates as the query string.
  \item The reverse-geocoded landmark string, if available.
\end{itemize}

WhatsApp-dominant set includes IN, BR, GB, NG, ZA, MX, ID, ES, IT, AR, CO, CL, PE, EG, SA, AE, MY, TR, and roughly 70 others. Because \code{wa.me} accepts only one recipient per URL, this path sends to a single contact at a time.

\subsection{SMS-Dominant Countries}

For countries dominated by SMS (US, CA, FR, DE, JP, KR, AU, and most of Northern Europe), the handler issues \code{window.location.href = 'sms:<contact1>,<contact2>,\dots?body=<encoded message>'}. The native SMS composer opens with all configured emergency contacts pre-populated as recipients.

\subsection{Event Telemetry}

After dispatch, the page emits a \code{roadsos:sos-sent} \code{CustomEvent} with details about the recipient and the medium chosen. \file{App.jsx} listens for this event and opens \file{DispatchScreen.jsx} for the post-send confirmation UI. The same event triggers a fire-and-forget \code{POST /dispatch} to the backend for telemetry.

% ──────────────────────────────────────────────────────────────────────
\section{Backend API Reference}
% ──────────────────────────────────────────────────────────────────────

\subsection{\texttt{GET /search}}

\begin{tabularx}{\textwidth}{l X}
\toprule
\textbf{Parameter} & \textbf{Description} \\
\midrule
\code{lat} & Float, WGS-84 latitude, range $[-90, 90]$. Required. \\
\code{lon} & Float, WGS-84 longitude, range $[-180, 180]$. Required. \\
\bottomrule
\end{tabularx}

Returns a JSON object with the shape given in Section~2.5. Rate limit: 30 req/min/IP.

\subsection{\texttt{POST /triage}}

\begin{lstlisting}[language=json,numbers=none,caption={Request body}]
{
  "contacts": [<list of contact objects from /search>],
  "injured": true,
  "blocking_traffic": false
}
\end{lstlisting}

Returns the same contact list, reordered, plus a \code{reason} string of one short sentence. Rate limit: 20 req/min/IP.

\subsection{\texttt{POST /dispatch}}

Telemetry sink. Logs the dispatch event server-side with request ID and returns \code{202 Accepted}. No body validation beyond JSON parsing.

\subsection{\texttt{GET /health}}

Returns \code{\{"status": "ok", "version": "1.x.x", "timestamp": "2026-05-16T..."\}}. Used by uptime monitors and by the frontend's \code{startBackendWarmup()} to wake the Render dyno on app load.

\subsection{\texttt{GET /offline-pack}}

Returns the bundled facility JSON for client-side warm caching. \code{Cache-Control: public, max-age=604800} (7 days).

% ──────────────────────────────────────────────────────────────────────
\section{Continuous Integration}
% ──────────────────────────────────────────────────────────────────────

Three GitHub Actions workflows are defined in \file{.github/workflows/}.

\subsection{Frontend CI (\texttt{frontend-ci.yml})}

\begin{itemize}[leftmargin=*, itemsep=2pt, topsep=2pt]
  \item \textbf{Triggers.} \code{push} and \code{pull\_request} against \code{main} when \code{frontend/**} is touched.
  \item \textbf{Matrix.} Node.js 20 and 22 in parallel.
  \item \textbf{Steps.} \code{npm ci --legacy-peer-deps} $\rightarrow$ \code{vite build} $\rightarrow$ \code{vitest run}.
  \item \textbf{Artifact.} The Node-20 job uploads \code{frontend/dist} as a build artifact (7-day retention).
\end{itemize}

\subsection{Backend CI (\texttt{backend-tests.yml})}

Three jobs run sequentially:

\begin{enumerate}[leftmargin=*, itemsep=2pt, topsep=2pt]
  \item \textbf{lint.} \code{ruff check} (rule set: E, W, F, I, UP, B, C4, SIM) followed by \code{ruff format --check}. Five rules are explicitly ignored: E501 (line length), E402 (intentional \code{load\_dotenv()} ordering), B008 (FastAPI \code{Depends()}), SIM108 (ternary readability), and SIM112 (the \code{Mapsplatformkey} environment variable name is fixed by Render configuration).
  \item \textbf{startup-check.} Single-step verification that \code{python -c "import main"} succeeds against \code{requirements.txt} with dummy API keys --- catches dependency-chain breakage before tests run.
  \item \textbf{test.} \code{pytest} on Python 3.11 and 3.12 in parallel. Depends on \code{lint} and \code{startup-check} passing.
\end{enumerate}

\subsection{PR Guard (\texttt{pr-guard.yml})}

Two jobs run on every pull request against \code{main}:

\begin{enumerate}[leftmargin=*, itemsep=2pt, topsep=2pt]
  \item \textbf{conflict-check.} Attempts a no-commit, no-fast-forward merge of the PR branch against \code{origin/main}. Fails if Git reports \code{CONFLICT}, surfacing the issue before review.
  \item \textbf{branch-staleness.} Counts commits the PR branch is behind \code{main} via \code{git rev-list --count HEAD..origin/main}. If the count exceeds 30, emits a warning recommending a rebase. This prevents the kind of large-scale merge conflict that arose during the merge of \code{arindam-branch-} into \code{main}.
\end{enumerate}

\subsection{Branch Protection (Applied at the Repository Level)}

The \code{main} branch enforces:

\begin{itemize}[leftmargin=*, itemsep=2pt, topsep=2pt]
  \item Required status checks: Build \& Test (Node 20), Build \& Test (Node 22), pytest (Python 3.11), pytest (Python 3.12), Lint (ruff), App startup check, and Merge conflict check.
  \item Strict mode: branches must be up to date with \code{main} before merging.
  \item No force pushes.
  \item No branch deletions.
\end{itemize}

\subsection{Test Coverage}

Backend tests are organised by service:

\begin{tabularx}{\textwidth}{l c X}
\toprule
\textbf{Test file} & \textbf{LOC} & \textbf{Scope} \\
\midrule
\code{test\_search\_service.py} & 63 & deduplication logic: invalid dicts, empty names, case-insensitive name match, phone equality \\
\code{test\_ai\_triage.py} & 90 & rule-based fallback, JSON schema validation, single- and two-leg priority orderings \\
\code{test\_overpass.py} & 207 & element parsing, haversine accuracy, proximity dedup, category classification \\
\code{test\_overpass\_parsing.py} & --- & edge-case element shapes (nwr without center, tags missing) \\
\code{test\_rate\_limiter.py} & --- & token-bucket refill, \code{X-Forwarded-For} extraction, retry-after header \\
\code{test\_cache.py} & --- & TTL expiry, LRU eviction, hit/miss statistics \\
\code{test\_geocoder.py} & --- & ISO country-code validation against the regex \\
\code{test\_hours\_parser.py} & --- & \code{opening\_hours} parsing: \code{24/7}, day ranges, time ranges, public holidays \\
\code{test\_phone\_utils.py} & --- & E.164 normalisation, \code{is\_dialable} (3--15 digits), \code{phones\_match} \\
\code{test\_middleware.py} & --- & request-ID stamping, log line format, error handler 503 response \\
\bottomrule
\end{tabularx}

% ──────────────────────────────────────────────────────────────────────
\section{Deployment}
% ──────────────────────────────────────────────────────────────────────

\subsection{Backend on Render}

\file{render.yaml} at the repository root defines a single web service:

\begin{lstlisting}[language={},caption={render.yaml}]
services:
  - type: web
    name: roadsos-api
    runtime: python
    rootDir: backend
    plan: free
    buildCommand: pip install -r requirements.txt
    startCommand: uvicorn main:app --host 0.0.0.0 --port $PORT
    envVars:
      - key: GEMINI_API_KEY        # required, free at aistudio.google.com
        sync: false
      - key: Mapsplatformkey       # optional, comma-separated for multi-key
        sync: false
      - key: PYTHON_VERSION
        value: 3.11.9
\end{lstlisting}

Free-tier specifics:

\begin{itemize}[leftmargin=*, itemsep=2pt, topsep=2pt]
  \item Dyno sleeps after 15 minutes idle; first request triggers a $\sim$30-second cold start.
  \item No persistent disk (acceptable; all data is upstream or in client browser).
  \item TLS terminated at edge; backend sees \code{X-Forwarded-For} only.
\end{itemize}

\code{startBackendWarmup()} in the frontend fires a \code{GET /health} immediately on app mount, paying the cold-start cost before the user runs their first search.

\subsection{Frontend on Vercel}

\file{vercel.json}:

\begin{lstlisting}[language=json,caption={vercel.json}]
{
  "version": 2,
  "buildCommand": "cd frontend && npm install && npm run build",
  "outputDirectory": "frontend/dist",
  "installCommand": "echo 'install handled in buildCommand'",
  "framework": null,
  "rewrites": [
    { "source": "/(.*)", "destination": "/index.html" }
  ]
}
\end{lstlisting}

The single rewrite ensures client-side routes resolve through \code{index.html} (SPA behaviour). Build environment variables:

\begin{itemize}[leftmargin=*, itemsep=2pt, topsep=2pt]
  \item \code{VITE\_API\_URL} --- e.g.\ \code{https://roadsos-api.onrender.com}
\end{itemize}

\subsection{Bundle Sizes}

\begin{tabularx}{\textwidth}{l r r}
\toprule
\textbf{Asset} & \textbf{Raw} & \textbf{Gzip} \\
\midrule
\code{dist/assets/index.js} & 610 KiB & 181 KiB \\
\code{dist/assets/index.css} & 80 KiB & 18 KiB \\
\code{dist/sw.js} & 25 KiB & 8 KiB \\
\code{dist/index.html} & 0.8 KiB & 0.4 KiB \\
Precache total & --- & $\sim$210 KiB \\
\bottomrule
\end{tabularx}

% ──────────────────────────────────────────────────────────────────────
\section{Observability and Logging}
% ──────────────────────────────────────────────────────────────────────

\subsection{Request Logging}

Every request emits one structured log line:

\begin{lstlisting}[numbers=none,frame=none]
2026-05-16T14:32:08Z [INFO] roadsos.request  GET /search 200 1247ms a8f3e9c2
\end{lstlisting}

Format: timestamp, level, logger, method, path, status, duration in milliseconds, first eight characters of the request ID.

\subsection{Service-Level Logs}

\begin{tabularx}{\textwidth}{l X}
\toprule
\textbf{Service} & \textbf{Notable log events} \\
\midrule
\code{overpass\_service} & retry attempts (WARN), cache hits (INFO), per-fetch summary (INFO) \\
\code{googleplaces\_service} & enrichment count per /search (INFO), key rotation (DEBUG) \\
\code{geocoder} & invalid country code from Nominatim (WARN) \\
\code{ai\_triage} & API call failure with exception type (WARN), JSON validation failure (WARN) \\
\code{rate\_limiter} & per-IP throttling events (INFO) \\
\bottomrule
\end{tabularx}

\subsection{Suppressed Noise}

\file{logging\_config.py} silences \code{httpx}, \code{httpcore}, and \code{anyio} below WARN level. These otherwise emit a connection-lifecycle line for every internal HTTP call.

% ──────────────────────────────────────────────────────────────────────
\section{Repository Layout}
% ──────────────────────────────────────────────────────────────────────

\begin{lstlisting}[language={},numbers=none,frame=none]
Roadproj/
|- backend/
|  |- main.py                    # FastAPI app entry, middleware wiring
|  |- middleware.py              # RequestID, RequestLog, ErrorHandling
|  |- logging_config.py          # log format + library noise suppression
|  |- requirements.txt           # runtime deps (FastAPI, httpx, phonenumbers, ...)
|  |- requirements-dev.txt       # adds pytest, pytest-asyncio, ruff
|  |- pytest.ini                 # asyncio_mode = auto, testpaths
|  |- ruff.toml                  # lint rule set + project ignores
|  |- services/
|  |  |- search_service.py       # /search orchestrator (4 phases)
|  |  |- overpass_service.py     # OSM Overpass client + parser
|  |  |- googleplaces_service.py # Google Places + multi-key rotation
|  |  |- geocoder.py             # Nominatim reverse geocode
|  |  |- ai_triage.py            # Gemini 2.0 Flash + rule fallback
|  |  |- triage_service.py       # /triage router
|  |  |- cache.py                # async TTL/LRU cache
|  |  |- rate_limiter.py         # per-IP token bucket
|  |  |- phone_utils.py          # phonenumbers wrapper
|  |  |- hours_parser.py         # OSM opening_hours parser
|  |  |- health_service.py       # /health endpoint
|  |  |- dispatch_service.py     # /dispatch telemetry endpoint
|  |  `- offline_service.py      # /offline-pack endpoint
|  `- tests/                     # 10 test files, pytest discoverable
|
|- frontend/
|  |- vite.config.js             # Vite + vite-plugin-pwa configuration
|  |- vitest.config.js           # test runner config
|  |- public/sw.js               # Workbox service worker source
|  `- src/
|     |- App.jsx                 # top-level orchestrator + state machine
|     |- main.tsx                # ReactDOM root, i18n init
|     |- final-design.css        # design system: colours, typography, layout
|     |- components/
|     |  |- MapHero.jsx          # hero section (map + dock + SOS)
|     |  |- RealMap.jsx          # Leaflet map + custom markers
|     |  |- SOSButton.jsx        # SOS dispatch with WhatsApp/SMS routing
|     |  |- LanguagePicker.jsx   # first-launch 43-language modal
|     |  |- MedicalIdModal.jsx   # blood type, allergies, contacts
|     |  |- CrashAlert.jsx       # post-crash confirmation UI
|     |  |- DispatchScreen.jsx   # post-SOS confirmation UI
|     |  |- TriageModal.jsx      # injured + blocking questions
|     |  |- CountryEmergency.jsx # 108/112/911 banner
|     |  `- ...
|     |- hooks/
|     |  |- useLocation.js       # GPS + motion + crash detection
|     |  `- useNetwork.js        # online/offline status
|     |- utils/
|     |  |- overpass.js          # /search client
|     |  |- googlePlaces.js      # /triage client
|     |  |- offlineDB.js         # localStorage cache
|     |  |- bundledFacilities.js # tier-3 fallback search
|     |  |- emergencyNumbers.js  # country -> dial codes
|     |  |- sosDispatch.js       # WhatsApp/SMS deeplink builder
|     |  |- medicalId.js         # on-device emergency-contact storage
|     |  |- plusCodes.js         # OLC encoder (pure JS, no deps)
|     |  `- backendWarmup.js     # Render cold-start mitigation
|     |- i18n/
|     |  |- index.js             # i18next init + RTL handler
|     |  |- locales.js           # locale metadata (native, English, dir)
|     |  |- countryLanguage.js   # ISO country -> language code
|     |  `- *.json               # 43 translation bundles
|     `- data/
|        `- bundled_facilities.json # 249 facilities, 196 countries
|
|- docs/
|  |- ARCHITECTURE.md            # architecture review + ADRs
|  `- TECHNICAL.tex              # this document
|
|- .github/workflows/
|  |- frontend-ci.yml            # build + test on Node 20, 22
|  |- backend-tests.yml          # lint + startup + pytest
|  `- pr-guard.yml               # merge conflict + staleness detection
|
|- render.yaml                   # Render service definition
|- vercel.json                   # Vercel build configuration
|- LICENSE                       # MIT
`- README.md                     # human-facing overview
\end{lstlisting}

% ──────────────────────────────────────────────────────────────────────
\section{Key Numbers}
% ──────────────────────────────────────────────────────────────────────

\begin{tabularx}{\textwidth}{X r}
\toprule
\textbf{Metric} & \textbf{Value} \\
\midrule
Countries with at least one bundled facility & 196 \\
Bundled facilities total & 249 \\
Languages shipped & 43 \\
Indian Schedule-VIII languages covered & 22 of 22 \\
RTL languages supported & 6 (Arabic, Persian, Hebrew, Urdu, Kashmiri, Sindhi) \\
OSM categories queried in parallel & 9 \\
Overpass mirrors with retry & 3 \\
Overpass timeout per attempt & 30 s \\
Default Overpass radius & 5{,}000 m (expands to 10{,}000 m if sparse) \\
Bundled-fallback radius & 80 km (expands to 600 km if sparse) \\
Cache TTL --- Overpass / Google / geocode & 1 h / 1 h / 24 h \\
Cache key granularity & $\sim$11 m (4 decimal places) \\
localStorage cache TTL & 24 h \\
localStorage cache grid & $\sim$1.1 km \\
Rate limit --- /search & 30 req/min/IP \\
Rate limit --- /triage & 20 req/min/IP \\
Phone enrichment cap & 6 Place Details calls per /search \\
Crash detection --- velocity collapse threshold & $\ge$25 km/h $\rightarrow$ $\le$5 km/h within 2 s \\
Crash detection --- accelerometer threshold & $\ge$3.5 G \\
Crash detection --- signal alignment window & 4 s \\
Crash detection --- post-alert cooldown & 12 s \\
AI model & \code{gemini-2.0-flash} (free tier: 60 RPM / 1500 RPD) \\
Frontend bundle (gzip) & 181 KiB JS + 18 KiB CSS \\
Service worker precache entries & 6 ($\sim$210 KiB) \\
\bottomrule
\end{tabularx}

\vspace{0.8cm}
\hrule
\vspace{0.4cm}
\small
\noindent
\textbf{Document version.} Generated from \code{main} at commit \texttt{f94dc2a}. Re-run after substantive code changes; this is a hand-maintained reference, not auto-generated.

\end{document}
